% ============================================================
% f2012_liaoning_science_q18_prism.tex
% 2012 辽宁卷（理科）第 18 题：直三棱柱与二面角
% 来源 PDF：普通高考/2012/2012辽宁理.pdf
% 原页：第 1 页（A3），印面 751，第 18 题题干右侧
% 直三棱柱 ABC-A'B'C'，M 为 A'B 中点，N 为 B'C' 中点。
% 虚线表示被遮挡的棱 AB、AC、AA'。
% ============================================================
\begin{tikzpicture}[>=Stealth, line width=0.7pt, font=\small,
  every node/.style={inner sep=1.5pt}]
  \coordinate (A)  at (1.00,1.60);
  \coordinate (B)  at (0.20,0.20);
  \coordinate (C)  at (3.20,0.80);
  \coordinate (Ap) at (1.00,3.40);
  \coordinate (Bp) at (0.20,2.00);
  \coordinate (Cp) at (3.20,2.60);
  \coordinate (M)  at (0.60,1.80);
  \coordinate (N)  at (1.70,2.30);

  % 被遮挡的棱
  \draw[dashed] (A) -- (B);
  \draw[dashed] (A) -- (C);
  \draw[dashed] (A) -- (Ap);

  % 可见棱
  \draw (B) -- (C);
  \draw (B) -- (Bp);
  \draw (C) -- (Cp);
  \draw (Ap) -- (Bp);
  \draw (Bp) -- (Cp);
  \draw (Cp) -- (Ap);
  \draw (Ap) -- (B);

  % 中点连线与二面角相关线段
  \draw (M) -- (N);
  \draw (Ap) -- (N);
  \draw (C) -- (M);
  \draw (C) -- (N);

  % 点与标签
  \fill (A) circle (1.1pt) node[below right] {$A$};
  \fill (B) circle (1.1pt) node[below] {$B$};
  \fill (C) circle (1.1pt) node[right] {$C$};
  \fill (Ap) circle (1.1pt) node[above] {$A'$};
  \fill (Bp) circle (1.1pt) node[left] {$B'$};
  \fill (Cp) circle (1.1pt) node[right] {$C'$};
  \fill (M) circle (1.0pt) node[below left] {$M$};
  \fill (N) circle (1.0pt) node[above right] {$N$};
\end{tikzpicture}
